% =========================================================
% ppt/src/minggu-10.tex
% Slide pertemuan 10
% Bab XI — Operasi Logika, Pergeseran, dan Rotasi
%
% Sub-CPMK 3.2 | CPMK03 | CPL03
% =========================================================

\documentclass[aspectratio=169,11pt]{beamer}
% =========================================================
% ppt/src/preamble-ppt.tex
% Tema, warna, paket, dan makro bersama seluruh deck slide.
%
% Dipakai oleh setiap berkas minggu-NN.tex dengan pola:
%   \documentclass[aspectratio=169,11pt]{beamer}
%   % =========================================================
% ppt/src/preamble-ppt.tex
% Tema, warna, paket, dan makro bersama seluruh deck slide.
%
% Dipakai oleh setiap berkas minggu-NN.tex dengan pola:
%   \documentclass[aspectratio=169,11pt]{beamer}
%   % =========================================================
% ppt/src/preamble-ppt.tex
% Tema, warna, paket, dan makro bersama seluruh deck slide.
%
% Dipakai oleh setiap berkas minggu-NN.tex dengan pola:
%   \documentclass[aspectratio=169,11pt]{beamer}
%   \input{preamble-ppt}
%
% Palet warna disalin dari buku_ajar/config/styles.tex agar
% identitas visual slide sama dengan buku ajar.
% =========================================================

% ---------------------------------------------------------
% 1. Paket
% ---------------------------------------------------------
\usepackage[T1]{fontenc}
\usepackage[utf8]{inputenc}
\usepackage[indonesian]{babel}
\usepackage{lmodern}
\usepackage{tikz}
\usetikzlibrary{arrows.meta, positioning, calc, shapes.geometric}
\usepackage{booktabs}
\usepackage{array}
\usepackage{amsmath}
\usepackage{amsfonts}
\usepackage{pifont}          % \ding untuk tanda checklist
\usepackage{listings}

% ---------------------------------------------------------
% 2. Palet warna OBE (identik dengan buku)
% ---------------------------------------------------------
\definecolor{obeBlue}{HTML}{1F3864}
\definecolor{obeBlueLight}{HTML}{DCE4F5}
\definecolor{obeTeal}{HTML}{0F6E6E}
\definecolor{obeTealLight}{HTML}{D9F0EE}
\definecolor{obeGreen}{HTML}{2E6B34}
\definecolor{obeGreenLight}{HTML}{DEEFDF}
\definecolor{obeAmber}{HTML}{8A5A00}
\definecolor{obeAmberLight}{HTML}{FDF1D6}
\definecolor{obeMaroon}{HTML}{7A2233}
\definecolor{obeMaroonLight}{HTML}{F7E1E5}
\definecolor{obeGray}{HTML}{4A4A4A}
\definecolor{obeGrayLight}{HTML}{F0F0F0}
\definecolor{codeBack}{HTML}{F7F7F4}
\definecolor{codeFrame}{HTML}{C9C9C0}

% ---------------------------------------------------------
% 3. Identitas mata kuliah (dipakai seluruh deck)
% ---------------------------------------------------------
\newcommand{\NamaMK}{Bahasa Pemrograman Tingkat Rendah}
\newcommand{\NamaMKShort}{Bhs.\ Pemrograman Tingkat Rendah}
\newcommand{\KodeMK}{MWCA010}
\newcommand{\Prodi}{Program Studi Teknik Informatika}
\newcommand{\Fakultas}{Fakultas Teknologi Informasi}
\newcommand{\PerguruanTinggi}{Universitas Bale Bandung}
\newcommand{\DosenPengampu}{Cecep Suwanda, S.Si., M.Kom.}
\newcommand{\SemesterTeks}{Semester Ganjil}

% \PertemuanTeks dan \JudulBab ditetapkan di masing-masing
% berkas minggu-NN.tex sebelum \begin{document}.
\providecommand{\PertemuanTeks}{Pertemuan}

% ---------------------------------------------------------
% 4. Tema beamer
% ---------------------------------------------------------
\usetheme{default}
\setbeamertemplate{navigation symbols}{}

\setbeamercolor{normal text}{fg=black!88}
\setbeamercolor{structure}{fg=obeBlue}
\setbeamercolor{alerted text}{fg=obeMaroon}

% --- judul slide: pita biru penuh lebar ---
\setbeamercolor{frametitle}{fg=white,bg=obeBlue}
\setbeamerfont{frametitle}{series=\bfseries,size=\large}
\setbeamertemplate{frametitle}{%
    \nointerlineskip%
    \begin{beamercolorbox}[%
        wd=\paperwidth,ht=2.9ex,dp=1.5ex,
        leftskip=0.4cm,rightskip=0.4cm]{frametitle}%
        \usebeamerfont{frametitle}\insertframetitle%
    \end{beamercolorbox}%
    \nointerlineskip%
    \hbox{\textcolor{obeTeal}{\rule{\paperwidth}{1.1pt}}}%
}

% --- kaki slide: identitas mata kuliah, pertemuan, nomor ---
\setbeamercolor{footline}{fg=white,bg=obeGray}
\setbeamerfont{footline}{size=\scriptsize,family=\sffamily}
\setbeamertemplate{footline}{%
    \leavevmode%
    \hbox{%
        \begin{beamercolorbox}[wd=0.42\paperwidth,ht=2.1ex,dp=1.1ex,
            leftskip=0.4cm]{footline}%
            \usebeamerfont{footline}%
            \KodeMK\ \textbullet\ \NamaMKShort%
        \end{beamercolorbox}%
        \begin{beamercolorbox}[wd=0.32\paperwidth,ht=2.1ex,dp=1.1ex,
            center]{footline}%
            \usebeamerfont{footline}\PertemuanTeks%
        \end{beamercolorbox}%
        \begin{beamercolorbox}[wd=0.26\paperwidth,ht=2.1ex,dp=1.1ex,
            rightskip=0.4cm]{footline}%
            \usebeamerfont{footline}%
            \insertframenumber\,/\,\inserttotalframenumber%
        \end{beamercolorbox}%
    }%
    \vskip0pt%
}

% --- daftar ---
\setbeamercolor{itemize item}{fg=obeTeal}
\setbeamercolor{itemize subitem}{fg=obeTeal}
\setbeamercolor{enumerate item}{fg=obeBlue}
\setbeamertemplate{itemize item}{\raise0.2ex\hbox{\tiny$\blacksquare$}}
\setbeamertemplate{itemize subitem}{\raise0.2ex\hbox{\scriptsize$\bullet$}}

% --- kotak (pengganti tcolorbox pada buku) ---
\setbeamercolor{block title}{bg=obeBlueLight,fg=obeBlue}
\setbeamercolor{block body}{bg=obeBlueLight!45,fg=black!88}
\setbeamercolor{block title alerted}{bg=obeMaroonLight,fg=obeMaroon}
\setbeamercolor{block body alerted}{bg=obeMaroonLight!45,fg=black!88}
\setbeamercolor{block title example}{bg=obeTealLight,fg=obeTeal}
\setbeamercolor{block body example}{bg=obeTealLight!45,fg=black!88}

% ---------------------------------------------------------
% 5. Makro notasi (sama persis dengan buku)
%
% Catatan: \note TIDAK didefinisikan di sini karena beamer
% sudah memakai \note untuk catatan pembicara. Untuk catatan
% di slide dipakai \catatan.
% ---------------------------------------------------------
\newcommand{\code}[1]{\texttt{#1}}
\newcommand{\asm}[1]{\texttt{#1}}
\newcommand{\reg}[1]{\texttt{#1}}
\newcommand{\flag}[1]{\texttt{#1}}
\newcommand{\dir}[1]{\texttt{#1}}
\newcommand{\file}[1]{\texttt{#1}}
\newcommand{\eng}[1]{\textit{#1}}

\newcommand{\hex}[1]{\texttt{#1h}}
\newcommand{\bin}[1]{\texttt{#1b}}
\newcommand{\dec}[1]{\texttt{#1}}

\newcommand{\important}[1]{\textbf{Penting:} #1}
\newcommand{\catatan}[1]{\textbf{Catatan:} #1}
\newcommand{\term}[1]{\textbf{#1}}

\newcommand{\CPL}[1]{\textbf{CPL#1}}
\newcommand{\CPMK}[1]{\textbf{CPMK#1}}
\newcommand{\SubCPMK}[1]{\textbf{Sub-CPMK #1}}

\newcommand{\checkbox}{\ding{113}\hspace{0.35em}}
\newcommand{\checked}{\ding{51}\hspace{0.35em}}
\newcommand{\crossed}{\ding{55}\hspace{0.35em}}

\newcommand{\flowarrow}{\ensuremath{\rightarrow}}

% ---------------------------------------------------------
% 6. Makro tata letak slide
% ---------------------------------------------------------

% Keterangan sumber materi (ketertelusuran ke buku)
\newcommand{\Sumber}[1]{%
    \vfill
    {\scriptsize\color{obeGray}Sumber: #1\par}%
}

% Frame judul deck: \JudulDeck{<judul bab>}{<pertemuan>}
\newcommand{\JudulDeck}[2]{%
\begin{frame}[plain]
    \vspace{0.25cm}
    {\color{obeBlue}\sffamily\bfseries\small \NamaMK\ (\KodeMK)\par}
    \vspace{0.12cm}
    {\color{obeGray}\sffamily\scriptsize
     \Prodi\ \textbullet\ \Fakultas\ \textbullet\ \PerguruanTinggi\par}
    \vfill
    {\color{obeBlue}\sffamily\LARGE\bfseries #1\par}
    \vspace{0.4cm}
    {\color{obeTeal}\sffamily\large #2\par}
    \vspace{0.45cm}
    {\color{obeTeal}\rule{\textwidth}{1pt}}
    \vfill
    {\color{obeGray}\sffamily\small Dosen Pengampu\par}
    \vspace{0.08cm}
    {\sffamily\bfseries \DosenPengampu\par}
    \vspace{0.2cm}
    {\color{obeGray}\sffamily\small \SemesterTeks\par}
    \vspace{0.25cm}
\end{frame}
}

% Frame pemisah bagian
\newcommand{\PemisahBagian}[1]{%
\begin{frame}[plain]
    \vfill
    \begin{center}
        {\color{obeTeal}\sffamily\small\bfseries BAGIAN\par}
        \vspace{0.35cm}
        {\color{obeBlue}\sffamily\Large\bfseries #1\par}
        \vspace{0.35cm}
        {\color{obeTeal}\rule{0.55\textwidth}{0.9pt}}
    \end{center}
    \vfill
\end{frame}
}

% ---------------------------------------------------------
% 7. Gaya penampil kode
% ---------------------------------------------------------
\input{listing-ppt}

%
% Palet warna disalin dari buku_ajar/config/styles.tex agar
% identitas visual slide sama dengan buku ajar.
% =========================================================

% ---------------------------------------------------------
% 1. Paket
% ---------------------------------------------------------
\usepackage[T1]{fontenc}
\usepackage[utf8]{inputenc}
\usepackage[indonesian]{babel}
\usepackage{lmodern}
\usepackage{tikz}
\usetikzlibrary{arrows.meta, positioning, calc, shapes.geometric}
\usepackage{booktabs}
\usepackage{array}
\usepackage{amsmath}
\usepackage{amsfonts}
\usepackage{pifont}          % \ding untuk tanda checklist
\usepackage{listings}

% ---------------------------------------------------------
% 2. Palet warna OBE (identik dengan buku)
% ---------------------------------------------------------
\definecolor{obeBlue}{HTML}{1F3864}
\definecolor{obeBlueLight}{HTML}{DCE4F5}
\definecolor{obeTeal}{HTML}{0F6E6E}
\definecolor{obeTealLight}{HTML}{D9F0EE}
\definecolor{obeGreen}{HTML}{2E6B34}
\definecolor{obeGreenLight}{HTML}{DEEFDF}
\definecolor{obeAmber}{HTML}{8A5A00}
\definecolor{obeAmberLight}{HTML}{FDF1D6}
\definecolor{obeMaroon}{HTML}{7A2233}
\definecolor{obeMaroonLight}{HTML}{F7E1E5}
\definecolor{obeGray}{HTML}{4A4A4A}
\definecolor{obeGrayLight}{HTML}{F0F0F0}
\definecolor{codeBack}{HTML}{F7F7F4}
\definecolor{codeFrame}{HTML}{C9C9C0}

% ---------------------------------------------------------
% 3. Identitas mata kuliah (dipakai seluruh deck)
% ---------------------------------------------------------
\newcommand{\NamaMK}{Bahasa Pemrograman Tingkat Rendah}
\newcommand{\NamaMKShort}{Bhs.\ Pemrograman Tingkat Rendah}
\newcommand{\KodeMK}{MWCA010}
\newcommand{\Prodi}{Program Studi Teknik Informatika}
\newcommand{\Fakultas}{Fakultas Teknologi Informasi}
\newcommand{\PerguruanTinggi}{Universitas Bale Bandung}
\newcommand{\DosenPengampu}{Cecep Suwanda, S.Si., M.Kom.}
\newcommand{\SemesterTeks}{Semester Ganjil}

% \PertemuanTeks dan \JudulBab ditetapkan di masing-masing
% berkas minggu-NN.tex sebelum \begin{document}.
\providecommand{\PertemuanTeks}{Pertemuan}

% ---------------------------------------------------------
% 4. Tema beamer
% ---------------------------------------------------------
\usetheme{default}
\setbeamertemplate{navigation symbols}{}

\setbeamercolor{normal text}{fg=black!88}
\setbeamercolor{structure}{fg=obeBlue}
\setbeamercolor{alerted text}{fg=obeMaroon}

% --- judul slide: pita biru penuh lebar ---
\setbeamercolor{frametitle}{fg=white,bg=obeBlue}
\setbeamerfont{frametitle}{series=\bfseries,size=\large}
\setbeamertemplate{frametitle}{%
    \nointerlineskip%
    \begin{beamercolorbox}[%
        wd=\paperwidth,ht=2.9ex,dp=1.5ex,
        leftskip=0.4cm,rightskip=0.4cm]{frametitle}%
        \usebeamerfont{frametitle}\insertframetitle%
    \end{beamercolorbox}%
    \nointerlineskip%
    \hbox{\textcolor{obeTeal}{\rule{\paperwidth}{1.1pt}}}%
}

% --- kaki slide: identitas mata kuliah, pertemuan, nomor ---
\setbeamercolor{footline}{fg=white,bg=obeGray}
\setbeamerfont{footline}{size=\scriptsize,family=\sffamily}
\setbeamertemplate{footline}{%
    \leavevmode%
    \hbox{%
        \begin{beamercolorbox}[wd=0.42\paperwidth,ht=2.1ex,dp=1.1ex,
            leftskip=0.4cm]{footline}%
            \usebeamerfont{footline}%
            \KodeMK\ \textbullet\ \NamaMKShort%
        \end{beamercolorbox}%
        \begin{beamercolorbox}[wd=0.32\paperwidth,ht=2.1ex,dp=1.1ex,
            center]{footline}%
            \usebeamerfont{footline}\PertemuanTeks%
        \end{beamercolorbox}%
        \begin{beamercolorbox}[wd=0.26\paperwidth,ht=2.1ex,dp=1.1ex,
            rightskip=0.4cm]{footline}%
            \usebeamerfont{footline}%
            \insertframenumber\,/\,\inserttotalframenumber%
        \end{beamercolorbox}%
    }%
    \vskip0pt%
}

% --- daftar ---
\setbeamercolor{itemize item}{fg=obeTeal}
\setbeamercolor{itemize subitem}{fg=obeTeal}
\setbeamercolor{enumerate item}{fg=obeBlue}
\setbeamertemplate{itemize item}{\raise0.2ex\hbox{\tiny$\blacksquare$}}
\setbeamertemplate{itemize subitem}{\raise0.2ex\hbox{\scriptsize$\bullet$}}

% --- kotak (pengganti tcolorbox pada buku) ---
\setbeamercolor{block title}{bg=obeBlueLight,fg=obeBlue}
\setbeamercolor{block body}{bg=obeBlueLight!45,fg=black!88}
\setbeamercolor{block title alerted}{bg=obeMaroonLight,fg=obeMaroon}
\setbeamercolor{block body alerted}{bg=obeMaroonLight!45,fg=black!88}
\setbeamercolor{block title example}{bg=obeTealLight,fg=obeTeal}
\setbeamercolor{block body example}{bg=obeTealLight!45,fg=black!88}

% ---------------------------------------------------------
% 5. Makro notasi (sama persis dengan buku)
%
% Catatan: \note TIDAK didefinisikan di sini karena beamer
% sudah memakai \note untuk catatan pembicara. Untuk catatan
% di slide dipakai \catatan.
% ---------------------------------------------------------
\newcommand{\code}[1]{\texttt{#1}}
\newcommand{\asm}[1]{\texttt{#1}}
\newcommand{\reg}[1]{\texttt{#1}}
\newcommand{\flag}[1]{\texttt{#1}}
\newcommand{\dir}[1]{\texttt{#1}}
\newcommand{\file}[1]{\texttt{#1}}
\newcommand{\eng}[1]{\textit{#1}}

\newcommand{\hex}[1]{\texttt{#1h}}
\newcommand{\bin}[1]{\texttt{#1b}}
\newcommand{\dec}[1]{\texttt{#1}}

\newcommand{\important}[1]{\textbf{Penting:} #1}
\newcommand{\catatan}[1]{\textbf{Catatan:} #1}
\newcommand{\term}[1]{\textbf{#1}}

\newcommand{\CPL}[1]{\textbf{CPL#1}}
\newcommand{\CPMK}[1]{\textbf{CPMK#1}}
\newcommand{\SubCPMK}[1]{\textbf{Sub-CPMK #1}}

\newcommand{\checkbox}{\ding{113}\hspace{0.35em}}
\newcommand{\checked}{\ding{51}\hspace{0.35em}}
\newcommand{\crossed}{\ding{55}\hspace{0.35em}}

\newcommand{\flowarrow}{\ensuremath{\rightarrow}}

% ---------------------------------------------------------
% 6. Makro tata letak slide
% ---------------------------------------------------------

% Keterangan sumber materi (ketertelusuran ke buku)
\newcommand{\Sumber}[1]{%
    \vfill
    {\scriptsize\color{obeGray}Sumber: #1\par}%
}

% Frame judul deck: \JudulDeck{<judul bab>}{<pertemuan>}
\newcommand{\JudulDeck}[2]{%
\begin{frame}[plain]
    \vspace{0.25cm}
    {\color{obeBlue}\sffamily\bfseries\small \NamaMK\ (\KodeMK)\par}
    \vspace{0.12cm}
    {\color{obeGray}\sffamily\scriptsize
     \Prodi\ \textbullet\ \Fakultas\ \textbullet\ \PerguruanTinggi\par}
    \vfill
    {\color{obeBlue}\sffamily\LARGE\bfseries #1\par}
    \vspace{0.4cm}
    {\color{obeTeal}\sffamily\large #2\par}
    \vspace{0.45cm}
    {\color{obeTeal}\rule{\textwidth}{1pt}}
    \vfill
    {\color{obeGray}\sffamily\small Dosen Pengampu\par}
    \vspace{0.08cm}
    {\sffamily\bfseries \DosenPengampu\par}
    \vspace{0.2cm}
    {\color{obeGray}\sffamily\small \SemesterTeks\par}
    \vspace{0.25cm}
\end{frame}
}

% Frame pemisah bagian
\newcommand{\PemisahBagian}[1]{%
\begin{frame}[plain]
    \vfill
    \begin{center}
        {\color{obeTeal}\sffamily\small\bfseries BAGIAN\par}
        \vspace{0.35cm}
        {\color{obeBlue}\sffamily\Large\bfseries #1\par}
        \vspace{0.35cm}
        {\color{obeTeal}\rule{0.55\textwidth}{0.9pt}}
    \end{center}
    \vfill
\end{frame}
}

% ---------------------------------------------------------
% 7. Gaya penampil kode
% ---------------------------------------------------------
% =========================================================
% ppt/src/listing-ppt.tex
% Gaya penampil kode untuk slide.
%
% Diturunkan dari buku_ajar/config/listings.tex, tetapi
% disesuaikan untuk slide: ukuran lebih kecil (karena frame
% 16:9 dan kode harus utuh terlihat) dan jarak lebih rapat.
% =========================================================

% ---------------------------------------------------------
% 1. Definisi bahasa MASM/TASM 8086 (salinan dari buku)
% ---------------------------------------------------------
\lstdefinelanguage{MASM}{
    morekeywords={
        MOV, XCHG, LEA, LDS, LES, PUSH, POP, PUSHA, POPA, PUSHF, POPF,
        LAHF, SAHF, XLAT, IN, OUT, CBW, CWD,
        ADD, ADC, SUB, SBB, INC, DEC, NEG, CMP, MUL, IMUL, DIV, IDIV,
        AAA, AAS, AAM, AAD, DAA, DAS,
        AND, OR, XOR, NOT, TEST, SHL, SHR, SAL, SAR, ROL, ROR, RCL, RCR,
        JMP, JE, JZ, JNE, JNZ, JG, JGE, JL, JLE, JA, JAE, JB, JBE,
        JC, JNC, JS, JNS, JO, JNO, JP, JPE, JNP, JPO, JCXZ,
        LOOP, LOOPE, LOOPZ, LOOPNE, LOOPNZ, CALL, RET, INT, IRET, INTO,
        HLT, NOP, CLI, STI, CLD, STD, CLC, STC, CMC, WAIT, ESC, LOCK,
        DB, DW, DD, DQ, DT, EQU, ORG, SEGMENT, ENDS, ASSUME, END,
        PROC, ENDP, MACRO, ENDM, LOCAL, PUBLIC, EXTRN, EXTERN,
        INCLUDE, INCLUDELIB, MODEL, STACK, CODE, DATA, CONST, ALIGN, EVEN,
        OFFSET, PTR, BYTE, WORD, DWORD, QWORD, TBYTE,
        NEAR, FAR, SHORT, LABEL, TYPE, SIZE, LENGTH, DUP, THIS,
        AX, BX, CX, DX, AL, AH, BL, BH, CL, CH, DL, DH,
        SP, BP, SI, DI, IP, CS, DS, SS, ES, FLAGS,
        CF, PF, AF, ZF, TF, IF, DF, OF
    },
    sensitive=false,
    morecomment=[l]{;},
    morestring=[b]',
    morestring=[b]"
}

% ---------------------------------------------------------
% 2. Gaya kode assembly untuk slide
% ---------------------------------------------------------
\lstdefinestyle{asmslide}{
    language=MASM,
    basicstyle=\ttfamily\scriptsize,
    keywordstyle=\color{obeBlue}\bfseries,
    commentstyle=\color{obeGreen}\itshape,
    stringstyle=\color{obeMaroon},
    numberstyle=\tiny\color{obeGray},
    numbers=left,
    numbersep=5pt,
    stepnumber=1,
    backgroundcolor=\color{codeBack},
    frame=single,
    rulecolor=\color{codeFrame},
    framesep=4pt,
    breaklines=true,
    breakatwhitespace=false,
    % CATATAN PENTING: opsi `postbreak` TIDAK boleh dipakai di sini.
    % Di beamer, `breaklines=true` bersama `postbreak=<kode>` membuat
    % listings gagal dirakit dengan galat
    %   "Missing number, treated as zero <to be read again> \kern".
    % Karena \lstset{style=asmslide} bersifat global, opsi postbreak
    % juga bocor ke gaya lain (mis. outslide/keluaran) dan merusaknya.
    % Di buku ajar (kelas `book`) justru sebaliknya: postbreak aman dan
    % dipakai. Jadi larangan ini khusus untuk slide.
    showstringspaces=false,
    tabsize=4,
    columns=fullflexible,
    keepspaces=true,
    upquote=true,
    xleftmargin=12pt,
    aboveskip=0.4em,
    belowskip=0.4em
}

% ---------------------------------------------------------
% 3. Gaya keluaran program / sesi debugger untuk slide
% ---------------------------------------------------------
\lstdefinestyle{outslide}{
    basicstyle=\ttfamily\scriptsize,
    keywordstyle=\color{black},
    commentstyle=\color{obeGray},
    stringstyle=\color{black},
    numbers=none,
    backgroundcolor=\color{obeGrayLight},
    frame=single,
    rulecolor=\color{codeFrame},
    framesep=4pt,
    breaklines=true,
    showstringspaces=false,
    columns=fullflexible,
    keepspaces=true,
    upquote=true,
    xleftmargin=4pt,
    aboveskip=0.4em,
    belowskip=0.4em
}

\lstset{style=asmslide}

% ---------------------------------------------------------
% 4. Environment ringkas untuk dipakai di dalam frame
%
% Frame yang memuat kedua environment ini WAJIB memakai
% opsi [fragile], karena isinya verbatim.
% ---------------------------------------------------------
\lstnewenvironment{kodeasm}[1][]{%
    \lstset{style=asmslide, #1}%
}{}

\lstnewenvironment{keluaran}[1][]{%
    \lstset{style=outslide, #1}%
}{}


%
% Palet warna disalin dari buku_ajar/config/styles.tex agar
% identitas visual slide sama dengan buku ajar.
% =========================================================

% ---------------------------------------------------------
% 1. Paket
% ---------------------------------------------------------
\usepackage[T1]{fontenc}
\usepackage[utf8]{inputenc}
\usepackage[indonesian]{babel}
\usepackage{lmodern}
\usepackage{tikz}
\usetikzlibrary{arrows.meta, positioning, calc, shapes.geometric}
\usepackage{booktabs}
\usepackage{array}
\usepackage{amsmath}
\usepackage{amsfonts}
\usepackage{pifont}          % \ding untuk tanda checklist
\usepackage{listings}

% ---------------------------------------------------------
% 2. Palet warna OBE (identik dengan buku)
% ---------------------------------------------------------
\definecolor{obeBlue}{HTML}{1F3864}
\definecolor{obeBlueLight}{HTML}{DCE4F5}
\definecolor{obeTeal}{HTML}{0F6E6E}
\definecolor{obeTealLight}{HTML}{D9F0EE}
\definecolor{obeGreen}{HTML}{2E6B34}
\definecolor{obeGreenLight}{HTML}{DEEFDF}
\definecolor{obeAmber}{HTML}{8A5A00}
\definecolor{obeAmberLight}{HTML}{FDF1D6}
\definecolor{obeMaroon}{HTML}{7A2233}
\definecolor{obeMaroonLight}{HTML}{F7E1E5}
\definecolor{obeGray}{HTML}{4A4A4A}
\definecolor{obeGrayLight}{HTML}{F0F0F0}
\definecolor{codeBack}{HTML}{F7F7F4}
\definecolor{codeFrame}{HTML}{C9C9C0}

% ---------------------------------------------------------
% 3. Identitas mata kuliah (dipakai seluruh deck)
% ---------------------------------------------------------
\newcommand{\NamaMK}{Bahasa Pemrograman Tingkat Rendah}
\newcommand{\NamaMKShort}{Bhs.\ Pemrograman Tingkat Rendah}
\newcommand{\KodeMK}{MWCA010}
\newcommand{\Prodi}{Program Studi Teknik Informatika}
\newcommand{\Fakultas}{Fakultas Teknologi Informasi}
\newcommand{\PerguruanTinggi}{Universitas Bale Bandung}
\newcommand{\DosenPengampu}{Cecep Suwanda, S.Si., M.Kom.}
\newcommand{\SemesterTeks}{Semester Ganjil}

% \PertemuanTeks dan \JudulBab ditetapkan di masing-masing
% berkas minggu-NN.tex sebelum \begin{document}.
\providecommand{\PertemuanTeks}{Pertemuan}

% ---------------------------------------------------------
% 4. Tema beamer
% ---------------------------------------------------------
\usetheme{default}
\setbeamertemplate{navigation symbols}{}

\setbeamercolor{normal text}{fg=black!88}
\setbeamercolor{structure}{fg=obeBlue}
\setbeamercolor{alerted text}{fg=obeMaroon}

% --- judul slide: pita biru penuh lebar ---
\setbeamercolor{frametitle}{fg=white,bg=obeBlue}
\setbeamerfont{frametitle}{series=\bfseries,size=\large}
\setbeamertemplate{frametitle}{%
    \nointerlineskip%
    \begin{beamercolorbox}[%
        wd=\paperwidth,ht=2.9ex,dp=1.5ex,
        leftskip=0.4cm,rightskip=0.4cm]{frametitle}%
        \usebeamerfont{frametitle}\insertframetitle%
    \end{beamercolorbox}%
    \nointerlineskip%
    \hbox{\textcolor{obeTeal}{\rule{\paperwidth}{1.1pt}}}%
}

% --- kaki slide: identitas mata kuliah, pertemuan, nomor ---
\setbeamercolor{footline}{fg=white,bg=obeGray}
\setbeamerfont{footline}{size=\scriptsize,family=\sffamily}
\setbeamertemplate{footline}{%
    \leavevmode%
    \hbox{%
        \begin{beamercolorbox}[wd=0.42\paperwidth,ht=2.1ex,dp=1.1ex,
            leftskip=0.4cm]{footline}%
            \usebeamerfont{footline}%
            \KodeMK\ \textbullet\ \NamaMKShort%
        \end{beamercolorbox}%
        \begin{beamercolorbox}[wd=0.32\paperwidth,ht=2.1ex,dp=1.1ex,
            center]{footline}%
            \usebeamerfont{footline}\PertemuanTeks%
        \end{beamercolorbox}%
        \begin{beamercolorbox}[wd=0.26\paperwidth,ht=2.1ex,dp=1.1ex,
            rightskip=0.4cm]{footline}%
            \usebeamerfont{footline}%
            \insertframenumber\,/\,\inserttotalframenumber%
        \end{beamercolorbox}%
    }%
    \vskip0pt%
}

% --- daftar ---
\setbeamercolor{itemize item}{fg=obeTeal}
\setbeamercolor{itemize subitem}{fg=obeTeal}
\setbeamercolor{enumerate item}{fg=obeBlue}
\setbeamertemplate{itemize item}{\raise0.2ex\hbox{\tiny$\blacksquare$}}
\setbeamertemplate{itemize subitem}{\raise0.2ex\hbox{\scriptsize$\bullet$}}

% --- kotak (pengganti tcolorbox pada buku) ---
\setbeamercolor{block title}{bg=obeBlueLight,fg=obeBlue}
\setbeamercolor{block body}{bg=obeBlueLight!45,fg=black!88}
\setbeamercolor{block title alerted}{bg=obeMaroonLight,fg=obeMaroon}
\setbeamercolor{block body alerted}{bg=obeMaroonLight!45,fg=black!88}
\setbeamercolor{block title example}{bg=obeTealLight,fg=obeTeal}
\setbeamercolor{block body example}{bg=obeTealLight!45,fg=black!88}

% ---------------------------------------------------------
% 5. Makro notasi (sama persis dengan buku)
%
% Catatan: \note TIDAK didefinisikan di sini karena beamer
% sudah memakai \note untuk catatan pembicara. Untuk catatan
% di slide dipakai \catatan.
% ---------------------------------------------------------
\newcommand{\code}[1]{\texttt{#1}}
\newcommand{\asm}[1]{\texttt{#1}}
\newcommand{\reg}[1]{\texttt{#1}}
\newcommand{\flag}[1]{\texttt{#1}}
\newcommand{\dir}[1]{\texttt{#1}}
\newcommand{\file}[1]{\texttt{#1}}
\newcommand{\eng}[1]{\textit{#1}}

\newcommand{\hex}[1]{\texttt{#1h}}
\newcommand{\bin}[1]{\texttt{#1b}}
\newcommand{\dec}[1]{\texttt{#1}}

\newcommand{\important}[1]{\textbf{Penting:} #1}
\newcommand{\catatan}[1]{\textbf{Catatan:} #1}
\newcommand{\term}[1]{\textbf{#1}}

\newcommand{\CPL}[1]{\textbf{CPL#1}}
\newcommand{\CPMK}[1]{\textbf{CPMK#1}}
\newcommand{\SubCPMK}[1]{\textbf{Sub-CPMK #1}}

\newcommand{\checkbox}{\ding{113}\hspace{0.35em}}
\newcommand{\checked}{\ding{51}\hspace{0.35em}}
\newcommand{\crossed}{\ding{55}\hspace{0.35em}}

\newcommand{\flowarrow}{\ensuremath{\rightarrow}}

% ---------------------------------------------------------
% 6. Makro tata letak slide
% ---------------------------------------------------------

% Keterangan sumber materi (ketertelusuran ke buku)
\newcommand{\Sumber}[1]{%
    \vfill
    {\scriptsize\color{obeGray}Sumber: #1\par}%
}

% Frame judul deck: \JudulDeck{<judul bab>}{<pertemuan>}
\newcommand{\JudulDeck}[2]{%
\begin{frame}[plain]
    \vspace{0.25cm}
    {\color{obeBlue}\sffamily\bfseries\small \NamaMK\ (\KodeMK)\par}
    \vspace{0.12cm}
    {\color{obeGray}\sffamily\scriptsize
     \Prodi\ \textbullet\ \Fakultas\ \textbullet\ \PerguruanTinggi\par}
    \vfill
    {\color{obeBlue}\sffamily\LARGE\bfseries #1\par}
    \vspace{0.4cm}
    {\color{obeTeal}\sffamily\large #2\par}
    \vspace{0.45cm}
    {\color{obeTeal}\rule{\textwidth}{1pt}}
    \vfill
    {\color{obeGray}\sffamily\small Dosen Pengampu\par}
    \vspace{0.08cm}
    {\sffamily\bfseries \DosenPengampu\par}
    \vspace{0.2cm}
    {\color{obeGray}\sffamily\small \SemesterTeks\par}
    \vspace{0.25cm}
\end{frame}
}

% Frame pemisah bagian
\newcommand{\PemisahBagian}[1]{%
\begin{frame}[plain]
    \vfill
    \begin{center}
        {\color{obeTeal}\sffamily\small\bfseries BAGIAN\par}
        \vspace{0.35cm}
        {\color{obeBlue}\sffamily\Large\bfseries #1\par}
        \vspace{0.35cm}
        {\color{obeTeal}\rule{0.55\textwidth}{0.9pt}}
    \end{center}
    \vfill
\end{frame}
}

% ---------------------------------------------------------
% 7. Gaya penampil kode
% ---------------------------------------------------------
% =========================================================
% ppt/src/listing-ppt.tex
% Gaya penampil kode untuk slide.
%
% Diturunkan dari buku_ajar/config/listings.tex, tetapi
% disesuaikan untuk slide: ukuran lebih kecil (karena frame
% 16:9 dan kode harus utuh terlihat) dan jarak lebih rapat.
% =========================================================

% ---------------------------------------------------------
% 1. Definisi bahasa MASM/TASM 8086 (salinan dari buku)
% ---------------------------------------------------------
\lstdefinelanguage{MASM}{
    morekeywords={
        MOV, XCHG, LEA, LDS, LES, PUSH, POP, PUSHA, POPA, PUSHF, POPF,
        LAHF, SAHF, XLAT, IN, OUT, CBW, CWD,
        ADD, ADC, SUB, SBB, INC, DEC, NEG, CMP, MUL, IMUL, DIV, IDIV,
        AAA, AAS, AAM, AAD, DAA, DAS,
        AND, OR, XOR, NOT, TEST, SHL, SHR, SAL, SAR, ROL, ROR, RCL, RCR,
        JMP, JE, JZ, JNE, JNZ, JG, JGE, JL, JLE, JA, JAE, JB, JBE,
        JC, JNC, JS, JNS, JO, JNO, JP, JPE, JNP, JPO, JCXZ,
        LOOP, LOOPE, LOOPZ, LOOPNE, LOOPNZ, CALL, RET, INT, IRET, INTO,
        HLT, NOP, CLI, STI, CLD, STD, CLC, STC, CMC, WAIT, ESC, LOCK,
        DB, DW, DD, DQ, DT, EQU, ORG, SEGMENT, ENDS, ASSUME, END,
        PROC, ENDP, MACRO, ENDM, LOCAL, PUBLIC, EXTRN, EXTERN,
        INCLUDE, INCLUDELIB, MODEL, STACK, CODE, DATA, CONST, ALIGN, EVEN,
        OFFSET, PTR, BYTE, WORD, DWORD, QWORD, TBYTE,
        NEAR, FAR, SHORT, LABEL, TYPE, SIZE, LENGTH, DUP, THIS,
        AX, BX, CX, DX, AL, AH, BL, BH, CL, CH, DL, DH,
        SP, BP, SI, DI, IP, CS, DS, SS, ES, FLAGS,
        CF, PF, AF, ZF, TF, IF, DF, OF
    },
    sensitive=false,
    morecomment=[l]{;},
    morestring=[b]',
    morestring=[b]"
}

% ---------------------------------------------------------
% 2. Gaya kode assembly untuk slide
% ---------------------------------------------------------
\lstdefinestyle{asmslide}{
    language=MASM,
    basicstyle=\ttfamily\scriptsize,
    keywordstyle=\color{obeBlue}\bfseries,
    commentstyle=\color{obeGreen}\itshape,
    stringstyle=\color{obeMaroon},
    numberstyle=\tiny\color{obeGray},
    numbers=left,
    numbersep=5pt,
    stepnumber=1,
    backgroundcolor=\color{codeBack},
    frame=single,
    rulecolor=\color{codeFrame},
    framesep=4pt,
    breaklines=true,
    breakatwhitespace=false,
    % CATATAN PENTING: opsi `postbreak` TIDAK boleh dipakai di sini.
    % Di beamer, `breaklines=true` bersama `postbreak=<kode>` membuat
    % listings gagal dirakit dengan galat
    %   "Missing number, treated as zero <to be read again> \kern".
    % Karena \lstset{style=asmslide} bersifat global, opsi postbreak
    % juga bocor ke gaya lain (mis. outslide/keluaran) dan merusaknya.
    % Di buku ajar (kelas `book`) justru sebaliknya: postbreak aman dan
    % dipakai. Jadi larangan ini khusus untuk slide.
    showstringspaces=false,
    tabsize=4,
    columns=fullflexible,
    keepspaces=true,
    upquote=true,
    xleftmargin=12pt,
    aboveskip=0.4em,
    belowskip=0.4em
}

% ---------------------------------------------------------
% 3. Gaya keluaran program / sesi debugger untuk slide
% ---------------------------------------------------------
\lstdefinestyle{outslide}{
    basicstyle=\ttfamily\scriptsize,
    keywordstyle=\color{black},
    commentstyle=\color{obeGray},
    stringstyle=\color{black},
    numbers=none,
    backgroundcolor=\color{obeGrayLight},
    frame=single,
    rulecolor=\color{codeFrame},
    framesep=4pt,
    breaklines=true,
    showstringspaces=false,
    columns=fullflexible,
    keepspaces=true,
    upquote=true,
    xleftmargin=4pt,
    aboveskip=0.4em,
    belowskip=0.4em
}

\lstset{style=asmslide}

% ---------------------------------------------------------
% 4. Environment ringkas untuk dipakai di dalam frame
%
% Frame yang memuat kedua environment ini WAJIB memakai
% opsi [fragile], karena isinya verbatim.
% ---------------------------------------------------------
\lstnewenvironment{kodeasm}[1][]{%
    \lstset{style=asmslide, #1}%
}{}

\lstnewenvironment{keluaran}[1][]{%
    \lstset{style=outslide, #1}%
}{}



\renewcommand{\PertemuanTeks}{Pertemuan 10}

\begin{document}

% =========================================================
% 1. JUDUL
% =========================================================
\JudulDeck{Operasi Logika, Pergeseran, dan Rotasi}{Pertemuan 10}

% =========================================================
% 2. CAPAIAN DAN TUJUAN
% =========================================================
\begin{frame}{Capaian dan Tujuan Pembelajaran}
\begin{block}{Sub-CPMK 3.2}
\small Menerapkan operasi logika \eng{bitwise} serta operasi pergeseran
dan rotasi (\asm{AND}, \asm{OR}, \asm{XOR}, \asm{NOT}, \asm{TEST},
\asm{SHL}, \asm{SHR}, \asm{SAL}, \asm{SAR}, \asm{ROL}, \asm{ROR})
untuk manipulasi bit.
\end{block}

\vspace{0.1em}
{\small\textbf{CPMK yang didukung:} \CPMK{03}\quad\textbullet\quad
\textbf{CPL yang didukung:} \CPL{03}\par}

\vspace{0.05em}
\begin{exampleblock}{Setelah pertemuan ini, mahasiswa mampu:}
\footnotesize
\begin{enumerate}
    \item menjelaskan tabel kebenaran \asm{AND}, \asm{OR}, \asm{XOR},
          dan \asm{NOT};
    \item menjelaskan pengaruh operasi logika terhadap \flag{CF},
          \flag{OF}, \flag{ZF}, \flag{SF}, dan \flag{PF};
    \item membedakan \asm{TEST} dan \asm{CMP} serta memilih yang tepat;
    \item menerapkan penyalakan, pemadaman, pembalikan, dan pengujian
          bit ke-$n$ pada sebuah byte;
    \item membedakan pergeseran logis, pergeseran aritmetika, dan
          rotasi beserta peran \flag{CF};
    \item menggantikan perkalian atau pembagian pangkat dua dengan
          pergeseran;
    \item menyusun program manipulasi bit: paritas,
          \eng{packing}/\eng{unpacking} BCD, dan \eng{XOR cipher}.
\end{enumerate}
\end{exampleblock}
\end{frame}

% =========================================================
% 3. PETA MATERI
% =========================================================
\begin{frame}{Peta Materi}
\small
\begin{enumerate}
    \item \textbf{Operasi Logika pada Level Bit}\\
          {\footnotesize tabel kebenaran; pengaruh terhadap \eng{flag};
          \eng{masking} dan manipulasi bit ke-$n$; \asm{TEST} vs
          \asm{CMP}}

    \vspace{0.25em}
    \item \textbf{Pergeseran dan Rotasi}\\
          {\footnotesize \asm{SHL}/\asm{SHR} dan \asm{SAL}/\asm{SAR};
          \asm{ROL}, \asm{ROR}, \asm{RCL}, \asm{RCR}; perkalian dan
          pembagian cepat}

    \vspace{0.25em}
    \item \textbf{Studi Kasus Manipulasi Bit}\\
          {\footnotesize pemeriksaan paritas; \eng{packing} dan
          \eng{unpacking} BCD; \eng{XOR cipher} sederhana}
\end{enumerate}

\vspace{0.35em}
\begin{block}{Praktik pada pertemuan ini}
\small Menelusuri \file{paritas.asm} dengan \eng{Single Step}, mencatat
\reg{AL}, \reg{BL}, dan \flag{CF}; dilanjutkan manipulasi bit ke-$n$
serta \eng{packing} BCD. Ditutup dengan \textbf{Kuis 3}.
\end{block}
\end{frame}

% =========================================================
% BAGIAN 1
% =========================================================
\PemisahBagian{Operasi Logika pada Level Bit}

% ---------------------------------------------------------
\begin{frame}{Operasi Logika pada Level Bit}
\small Berbeda dari \asm{ADD} dan \asm{SUB} yang menjalarkan
\eng{carry} antarbit, operasi logika memperlakukan setiap posisi bit
sebagai satuan mandiri yang tidak memengaruhi tetangganya.

\vspace{0.3em}
\begin{center}
\begin{tabular}{@{}ccccc@{}}
    \toprule
    \textbf{Bit $A$} & \textbf{Bit $B$} & \textbf{\asm{AND}} &
    \textbf{\asm{OR}} & \textbf{\asm{XOR}} \\
    \midrule
    0 & 0 & 0 & 0 & 0 \\
    0 & 1 & 0 & 1 & 1 \\
    1 & 0 & 0 & 1 & 1 \\
    1 & 1 & 1 & 1 & 0 \\
    \bottomrule
\end{tabular}
\end{center}

\vspace{0.3em}
{\footnotesize Hasilnya hanya bergantung pada pasangan bit yang
bersesuaian; keempat instruksi bekerja pada operand 8 bit maupun
16 bit.\par}

\Sumber{Bab XI, Tabel Kebenaran dan Empat Instruksi Dasar}
\end{frame}

% ---------------------------------------------------------
\begin{frame}[fragile]{Sifat Empat Instruksi Logika}
\small
\begin{columns}[t]
\begin{column}{0.47\textwidth}
\begin{itemize}
    \item \asm{AND}: 1 hanya bila kedua bit 1 --- mempertahankan bit
          tertentu, memadamkan sisanya.
    \item \asm{OR}: 1 bila salah satu bit 1 --- menyalakan bit tanpa
          mengubah bit lain.
    \item \asm{XOR}: 1 bila kedua bit berbeda --- membalik bit.
    \item \asm{NOT}: membalik seluruh bit dan \emph{tidak}
          menyentuh \eng{flag}.
\end{itemize}
\end{column}
\begin{column}{0.49\textwidth}
\begin{kodeasm}[numbers=none]
MOV AL, 10100101b   ; AL = A5h
AND AL, 00001111b   ; AL = 05h
MOV BL, 11000011b   ; BL = C3h
OR  AL, BL          ; AL = C7h
MOV CL, 10100101b   ; CL = A5h
XOR CL, 11111111b   ; CL = 5Ah
NOT CL              ; CL = A5h
\end{kodeasm}
\end{column}
\end{columns}
\Sumber{Bab XI, Tabel Kebenaran dan Empat Instruksi Dasar}
\end{frame}

% ---------------------------------------------------------
\begin{frame}{Pengaruh Operasi Logika terhadap \eng{Flag}}
\footnotesize
\begin{center}
\begin{tabular}{@{}p{1.5cm}p{5.0cm}p{6.4cm}@{}}
    \toprule
    \textbf{Flag} & \textbf{Nilai setelah operasi logika} &
    \textbf{Alasan} \\
    \midrule
    \flag{CF} & Selalu 0 & Tidak ada \eng{carry} yang keluar dari bit
    paling kiri \\
    \flag{OF} & Selalu 0 & Operasi \eng{bitwise} tidak menghasilkan
    luapan bertanda \\
    \flag{ZF} & 1 bila seluruh bit hasil bernilai 0 &
    Memeriksa apakah suatu bidang bit kosong \\
    \flag{SF} & Menyalin bit paling kiri hasil & Hasil dianggap negatif
    bila ditafsirkan bertanda \\
    \flag{PF} & 1 bila cacah bit 1 pada byte rendah genap & Dipakai
    pada pemeriksaan paritas \\
    \flag{AF} & Tidak terdefinisi & Jangan diandalkan setelah operasi
    logika \\
    \bottomrule
\end{tabular}
\end{center}

\vspace{0.25em}
{\footnotesize\asm{NOT} satu-satunya pengecualian: tidak menyentuh
\eng{flag} sama sekali, sehingga \flag{ZF}, \flag{SF}, dan \flag{PF}
masih bernilai seperti sebelumnya.\par}
\Sumber{Bab XI, Pengaruh Operasi Logika terhadap Flag}
\end{frame}

% ---------------------------------------------------------
\begin{frame}[fragile]{Aturan Penulisan Operand}
\small Aturan legalitas operand sama dengan \asm{MOV} dan \asm{ADD}:
operand tujuan boleh register atau memori, operand sumber boleh
register, memori, atau nilai segera.

\vspace{0.3em}
\begin{kodeasm}[numbers=none]
MOV AL, nilai           ; DB + AL -> ukuran sudah sama
MOV AX, WORD PTR total  ; DW + AX -> WORD PTR menyesuaikan ukuran
AND AL, BYTE PTR nilai  ; AL (8 bit) dengan operand memori
ADD BL, CL              ; register dengan register
\end{kodeasm}

\vspace{0.3em}
\begin{alertblock}{Dua larangan yang tetap berlaku}
\footnotesize Memori ke memori tidak diizinkan sehingga perlu register
perantara, dan kedua operand wajib berukuran sama.
\end{alertblock}
\Sumber{Bab XI, Aturan Penulisan Operand}
\end{frame}

% ---------------------------------------------------------
\begin{frame}[fragile]{Manipulasi Bit ke-$n$ dengan \eng{Mask}}
\small Misalkan \reg{AL} berisi \hex{A5} = \bin{1010 0101} dan bit yang
disasar bit ke-3 (\hex{08}). \eng{Mask} adalah $2^n$: satu bit menyala
pada posisi yang dituju.

\vspace{0.2em}
{\footnotesize
\begin{center}
\begin{tabular}{@{}p{2.6cm}p{2.2cm}p{3.2cm}p{4.4cm}@{}}
    \toprule
    \textbf{Kehendak} & \textbf{Mask} & \textbf{Instruksi} &
    \textbf{Hasil pada \hex{A5}, bit 3} \\
    \midrule
    Menyalakan bit & mask & \asm{OR AL, mask} &
    \bin{1010 1101} = \hex{AD} \\
    Memadamkan bit & \asm{NOT mask} & \asm{AND AL, mask} &
    \bin{1010 0101} = \hex{A5} \\
    Membalik bit & mask & \asm{XOR AL, mask} &
    \bin{1010 1101} = \hex{AD} \\
    Menguji bit & mask & \asm{TEST AL, mask} &
    \flag{ZF} = 0 berarti bit 3 menyala \\
    \bottomrule
\end{tabular}
\end{center}
\par}

\vspace{0.2em}
\begin{kodeasm}[numbers=none]
OR   AL, 08h        ; nyalakan bit 3      (A5h -> ADh)
AND  AL, 0F7h       ; padamkan bit 3      (F7h = NOT 08h)
XOR  AL, 08h        ; balikkan bit 3
TEST AL, 08h        ; uji bit 3 tanpa mengubah AL
\end{kodeasm}
\Sumber{Bab XI, Penerapan Praktis: Manipulasi Bit Tertentu}
\end{frame}

% ---------------------------------------------------------
\begin{frame}{Memadamkan Bit: Kesalahan yang Sering Terjadi}
\begin{alertblock}{Nilai mask sering ditulis terbalik}
\small Untuk memadamkan bit ke-3 yang benar adalah
\asm{AND AL, 0F7h} --- bit ke-3 pada mask bernilai 0. Penulisan
\asm{AND AL, 08h} justru mempertahankan bit itu dan menghapus seluruh
bit lain.
\end{alertblock}

\vspace{0.4em}
\begin{block}{Mask penghapus dihitung saat berjalan}
\small Mask disimpan ke register, dibalik dengan \asm{NOT}, lalu dipakai
pada \asm{AND}. Karena \asm{NOT} tidak menyentuh \eng{flag}, keputusan
lompatan sebelumnya tetap sah.
\end{block}

\vspace{0.4em}
{\footnotesize\color{obeGray}Kesalahan kedua: menganggap \asm{AND}
menyetel \flag{CF}. Operasi logika selalu mengosongkan \flag{CF} dan
\flag{OF}.\par}
\Sumber{Bab XI, Penerapan Praktis: Manipulasi Bit Tertentu}
\end{frame}

% ---------------------------------------------------------
\begin{frame}{\asm{TEST} dan \asm{CMP}: Dua Instruksi Penguji}
\footnotesize Keduanya melakukan operasi tanpa menyimpan hasilnya;
yang tertinggal hanya jejak pada \eng{flag}.

\vspace{0.2em}
\begin{center}
\begin{tabular}{@{}p{2.1cm}p{3.5cm}p{7.1cm}@{}}
    \toprule
    \textbf{Aspek} & \asm{TEST} & \asm{CMP} \\
    \midrule
    Operasi internal & \asm{AND} tanpa menyimpan hasil &
    \asm{SUB} tanpa menyimpan hasil \\
    \flag{CF} dan \flag{OF} & Selalu 0 & Sesuai hasil pengurangan \\
    Pemakaian utama & Memeriksa keadaan bit tertentu &
    Membandingkan besar dua bilangan secara utuh \\
    Lompatan lazim & \asm{JZ}/\asm{JNZ}, \asm{JP}/\asm{JNP} &
    \asm{JE}/\asm{JNE}, \asm{JA}/\asm{JB}, \asm{JG}/\asm{JL} \\
    Contoh & \asm{TEST AL, 01h} lalu \asm{JNZ} &
    \asm{CMP AL, BL} lalu \asm{JA} \\
    \bottomrule
\end{tabular}
\end{center}

\vspace{0.25em}
{\footnotesize Kebiasaan khusus: \asm{TEST AL, AL} tidak mengubah isi
\reg{AL}, tetapi memperbarui \flag{ZF}, \flag{SF}, dan \flag{PF};
pola ini menggantikan \asm{CMP AL, 0} dengan kode dua byte lebih
pendek.\par}
\Sumber{Bab XI, TEST dan CMP: Dua Instruksi Penguji}
\end{frame}

% ---------------------------------------------------------
\begin{frame}{Rangkuman Bagian 1}
\small
\begin{itemize}
    \item Operasi logika memandang setiap bit sebagai satuan mandiri;
          aturannya mengikuti gerbang logika dasar dan tidak
          menjalarkan \eng{carry}.

    \vspace{0.22em}
    \item \asm{AND} untuk \eng{masking}, \asm{OR} untuk menyalakan,
          \asm{XOR} untuk membalik; \asm{NOT} membalik seluruh bit dan
          tidak menyentuh \eng{flag}.

    \vspace{0.22em}
    \item \flag{CF} dan \flag{OF} selalu 0 setelah \asm{AND},
          \asm{OR}, \asm{XOR}, dan \asm{TEST}, sedangkan \flag{ZF},
          \flag{SF}, dan \flag{PF} mengikuti hasil --- sumber
          kesalahan bila \flag{CF} dipercaya.

    \vspace{0.22em}
    \item Bit ke-$n$ disentuh dengan mask $2^n$: nyalakan dengan
          \asm{OR}, padamkan dengan \asm{AND} memakai komplemen mask,
          balikkan dengan \asm{XOR}, uji dengan \asm{TEST}.

    \vspace{0.22em}
    \item \asm{TEST} = \asm{AND} tanpa menyimpan hasil (memeriksa bit);
          \asm{CMP} = \asm{SUB} tanpa menyimpan hasil (membandingkan
          nilai utuh).
\end{itemize}
\end{frame}

% =========================================================
% BAGIAN 2
% =========================================================
\PemisahBagian{Pergeseran dan Rotasi}

% ---------------------------------------------------------
\begin{frame}[fragile]{Pergeseran Logis: \asm{SHL} dan \asm{SHR}}
\small Operasi logika bekerja pada posisi bit yang tetap; pergeseran
justru memindahkan isi bit ke posisi lain. Bit yang terdesak keluar
ditampung pada \flag{CF}.

\vspace{0.25em}
\begin{kodeasm}[numbers=none]
SHL AL, 1           ; geser AL satu bit ke kiri, bit kanan diisi 0
SHR AX, 1           ; geser AX satu bit ke kanan, bit kiri diisi 0
MOV CL, 4           ; jumlah pergeseran disimpan pada CL
SHL BX, CL          ; geser BX empat bit ke kiri
\end{kodeasm}

\vspace{0.25em}
\begin{block}{Jumlah pergeseran}
\small Ditulis sebagai angka 1 atau disimpan pada \reg{CL}. Pada 8086
nilai \reg{CL} dipakai apa adanya (sampai 255), sedangkan pada 80286 ke
atas hanya lima bit terendahnya (batas 31).
\end{block}
\Sumber{Bab XI, Pergeseran Logis: SHL dan SHR}
\end{frame}

% ---------------------------------------------------------
\begin{frame}{Pergeseran Aritmetika: \asm{SAL} dan \asm{SAR}}
\small
\begin{itemize}
    \item \asm{SAL} (\eng{Shift Arithmetic Left}) secara operasi
          \emph{identik} dengan \asm{SHL} dan menghasilkan kode mesin
          yang sama; namanya hanya menegaskan nilai dianggap bertanda.

    \vspace{0.3em}
    \item \asm{SAR} berbeda secara hakiki: bit yang masuk dari sisi
          kiri bukan nol, melainkan \emph{salinan bit tanda}, sehingga
          tanda bilangan terpelihara.
\end{itemize}

\vspace{0.4em}
\begin{alertblock}{Mengapa \asm{SAR} diperlukan}
\small Pada \bin{1111 1111} ($-1$ dalam komplemen dua), \asm{SHR} satu
bit memberi \bin{0111 1111} ($+127$) karena bit tanda terisi nol.
Bagi bilangan bertanda negatif, hanya \asm{SAR} yang benar.
\end{alertblock}
\Sumber{Bab XI, Pergeseran Aritmetika: SAL dan SAR}
\end{frame}

% ---------------------------------------------------------
\begin{frame}{Perbedaan Pergeseran pada Byte \hex{96}}
\begin{center}
\begin{tikzpicture}[
    sel/.style={draw=obeGray, line width=0.5pt, minimum width=0.52cm,
        minimum height=0.52cm, inner sep=0pt, font=\ttfamily\scriptsize},
    keluar/.style={sel, fill=obeMaroonLight, draw=obeMaroon},
    masuk/.style={sel, fill=obeAmberLight, draw=obeAmber, dashed},
    baris/.style={font=\footnotesize\bfseries, anchor=east, text=obeBlue},
    ket/.style={font=\scriptsize, text=obeGray, anchor=west,
        align=left, text width=3.3cm}
]
    \foreach \i/\v in {0/1,1/0,2/0,3/1,4/0,5/1,6/1,7/0}
        \node[sel] at (\i*0.55, 0) {\v};
    \node[baris] at (-1.25, 0) {Byte awal};
    \node[ket] at (5.0, 0) {\hex{96} = \bin{1001 0110} ($-106$ bertanda)};

    \foreach \i/\v in {0/0,1/0,2/1,3/0,4/1,5/1,6/0,7/0}
        \node[sel] at (\i*0.55, -0.95) {\v};
    \node[keluar] at (-0.75, -0.95) {1};
    \node[masuk] at (4.6, -0.95) {0};
    \node[baris] at (-1.25, -0.95) {SHL 1 / SAL 1};
    \node[ket] at (5.0, -0.95) {bit 7 keluar, hasil $44$, \flag{CF} = 1};

    \foreach \i/\v in {0/0,1/1,2/0,3/0,4/1,5/0,6/1,7/1}
        \node[sel] at (\i*0.55, -1.9) {\v};
    \node[masuk] at (-0.75, -1.9) {0};
    \node[keluar] at (4.6, -1.9) {0};
    \node[baris] at (-1.25, -1.9) {SHR 1};
    \node[ket] at (5.0, -1.9) {0 masuk dari kiri, hasil $+75$,
        \flag{CF} = 0};

    \foreach \i/\v in {0/1,1/1,2/0,3/0,4/1,5/0,6/1,7/1}
        \node[sel] at (\i*0.55, -2.85) {\v};
    \node[masuk] at (-0.75, -2.85) {1};
    \node[keluar] at (4.6, -2.85) {0};
    \node[baris] at (-1.25, -2.85) {SAR 1};
    \node[ket] at (5.0, -2.85) {bit tanda disalin masuk; hasil $-53$};

    \node[keluar] at (-0.75, -3.65) {};
    \node[font=\scriptsize, text=obeGray, anchor=west] at (-0.35, -3.65)
        {keluar ke \flag{CF}};
    \node[masuk] at (2.4, -3.65) {};
    \node[font=\scriptsize, text=obeGray, anchor=west] at (2.8, -3.65)
        {masuk dari sisi kosong};
\end{tikzpicture}
\end{center}

\vspace{0.1em}
{\footnotesize\color{obeGray}\asm{SHR} memberi \bin{0100 1011} ($+75$)
karena bit tanda dipaksa nol, sedangkan \asm{SAR} memberi
\bin{1100 1011} ($-53$) karena bit tanda disalin ulang.\par}
\Sumber{Bab XI, Pergeseran Aritmetika: SAL dan SAR}
\end{frame}

% ---------------------------------------------------------
\begin{frame}{Ringkasan Kedelapan Instruksi Geser dan Rotasi}
\scriptsize
\begin{center}
\begin{tabular}{@{}p{1.3cm}p{1.1cm}p{3.2cm}p{6.9cm}@{}}
    \toprule
    \textbf{Instruksi} & \textbf{Arah} & \textbf{Bit yang masuk} &
    \textbf{Flag terpengaruh} \\
    \midrule
    \asm{SHL} & kiri & 0 di bit paling kanan &
    \flag{CF} menerima bit paling kiri; \flag{SF}, \flag{ZF},
    \flag{PF} ikut hasil \\
    \asm{SAL} & kiri & 0 di bit paling kanan &
    sama dengan \asm{SHL} (kode mesin identik) \\
    \asm{SHR} & kanan & 0 di bit paling kiri &
    \flag{CF} menerima bit paling kanan; \flag{SF}, \flag{ZF},
    \flag{PF} ikut hasil \\
    \asm{SAR} & kanan & salinan bit tanda &
    \flag{CF} menerima bit paling kanan; \flag{SF} tetap sesuai tanda \\
    \asm{ROL} & kiri & bit keluar dari kiri &
    \flag{CF} menerima bit berputar; \flag{OF} benar bila jumlah 1 \\
    \asm{ROR} & kanan & bit keluar dari kanan &
    \flag{CF} menerima bit berputar; \flag{OF} benar bila jumlah 1 \\
    \asm{RCL} & kiri & \flag{CF} lama &
    \flag{CF} menerima bit paling kiri, lalu ikut berputar \\
    \asm{RCR} & kanan & \flag{CF} lama &
    \flag{CF} menerima bit paling kanan, lalu ikut berputar \\
    \bottomrule
\end{tabular}
\end{center}

\vspace{0.1em}
{\scriptsize\color{obeGray}Pada 8086 \flag{OF} hanya terdefinisi benar
bila jumlah pergeseran satu; \flag{CF} selalu berisi bit terakhir yang
keluar, berapa pun jumlah pergeserannya.\par}
\Sumber{Bab XI, Ringkasan Kedelapan Instruksi Geser dan Rotasi}
\end{frame}

% ---------------------------------------------------------
\begin{frame}[fragile]{Rotasi: \asm{ROL} dan \asm{ROR}}
\small Rotasi memindahkan bit secara melingkar: bit yang terdesak keluar
dari satu sisi masuk kembali ke sisi lain, sehingga tidak ada informasi
yang hilang. Bit yang berputar selalu melewati \flag{CF}.

\vspace{0.3em}
\begin{kodeasm}[numbers=none]
MOV AL, 10000001b   ; AL = 81h
ROL AL, 1           ; AL = 00000011b, CF = 1 (bit 7 pindah ke bawah)

MOV AL, 10000001b   ; AL = 81h
ROR AL, 1           ; AL = 11000000b, CF = 1 (bit 0 pindah ke atas)
\end{kodeasm}

\vspace{0.3em}
{\footnotesize\color{obeGray}\flag{CF} merefleksikan bit terakhir yang
melewatinya, bukan bit yang hilang seperti pada pergeseran.\par}
\Sumber{Bab XI, Rotasi: ROL, ROR, RCL, dan RCR}
\end{frame}

% ---------------------------------------------------------
\begin{frame}[fragile]{\asm{RCL}, \asm{RCR}, dan Pergeseran 32 Bit}
\small Pada \asm{RCL}, \flag{CF} menjadi bit tambahan di sebelah kiri
operand: bit paling kiri masuk ke \flag{CF}, sedangkan nilai \flag{CF}
yang lama masuk ke bit paling kanan. \asm{RCR} berlaku kebalikannya.

\vspace{0.2em}
\begin{kodeasm}[numbers=none]
CLC                 ; CF = 0
MOV AL, 10000001b   ; AL = 81h
RCL AL, 1           ; AL = 00000010b, CF = 1 (CF lama masuk ke bit 0)

; menggeser bilangan 32 bit (DX:AX) satu bit ke kanan
SAR DX, 1           ; word tinggi digeser, bit tanda disalin
RCR AX, 1           ; bit yang keluar dari DX masuk ke AX
\end{kodeasm}

\vspace{0.2em}
{\footnotesize Pasangan \asm{SAR} dan \asm{RCR} memindahkan bit dari
word tinggi ke word rendah sehingga bilangan 32 bit dapat digeser
sebagai satu kesatuan.\par}
\Sumber{Bab XI, Rotasi: ROL, ROR, RCL, dan RCR}
\end{frame}

% ---------------------------------------------------------
\begin{frame}[fragile]{Perkalian dan Pembagian Cepat dengan Pergeseran}
\small Pergeseran satu bit setara dengan mengalikan atau membagi dengan
dua, sehingga pangkat dua tidak perlu memakai \asm{MUL} atau \asm{DIV}.

\vspace{0.2em}
\begin{kodeasm}[numbers=none]
MOV AX, 5
MOV CL, 4           ; 16 = 2 pangkat 4
SHL AX, CL          ; AX = 80  (5 x 16)

MOV AX, 5
MOV BX, AX          ; BX = 5, salinan untuk 5 x 8
SHL AX, 1           ; AX = 10  (5 x 2)
MOV CL, 3
SHL BX, CL          ; BX = 40  (5 x 8)
ADD AX, BX          ; AX = 50  (10 + 40)
\end{kodeasm}

\vspace{0.15em}
{\footnotesize\color{obeGray}Bentuk terakhir memakai $5 \times 10 =
(5 \times 2) + (5 \times 8)$; \asm{SHL} dan \asm{ADD} hanya perlu
beberapa siklus, sedangkan \asm{MUL} pada 8086 lebih dari seratus
siklus.\par}
\Sumber{Bab XI, Perkalian dan Pembagian Cepat dengan Pergeseran}
\end{frame}

% ---------------------------------------------------------
\begin{frame}{Peringatan: Membagi Bilangan Bertanda}
\begin{alertblock}{Jangan memakai \asm{SHR} untuk bilangan bertanda}
\small \asm{SHR AX, 1} mengosongkan bit tanda, sehingga $-1$
(\hex{FFFF}) menjadi \hex{7FFF} ($+32767$). Gunakan \asm{SAR}, bukan
\asm{SHR}.
\end{alertblock}

\vspace{0.5em}
\begin{block}{Namun \asm{SAR} tidak membulatkan ke arah nol}
\small \asm{SAR AX, 1} pada $-1$ menghasilkan $-1$, bukan $0$.
Instruksi \asm{IDIV} tetap diperlukan bila hasil pembagian harus
dibulatkan ke arah nol.
\end{block}

\vspace{0.5em}
{\footnotesize\color{obeGray}Pergeseran kanan logis hanya sah untuk
pembagian bilangan tak bertanda.\par}
\Sumber{Bab XI, Perkalian dan Pembagian Cepat dengan Pergeseran}
\end{frame}

% ---------------------------------------------------------
\begin{frame}{Rangkuman Bagian 2}
\small
\begin{itemize}
    \item \asm{SHL} dan \asm{SAL} identik (kode mesin sama) dan
          menyisipkan 0 dari kanan; \asm{SHR} menyisipkan 0 dari kiri;
          \asm{SAR} menyisipkan salinan bit tanda.

    \vspace{0.22em}
    \item Jumlah pergeseran ditulis 1 atau disimpan pada \reg{CL}.
          \flag{CF} selalu berisi bit terakhir yang keluar; \flag{OF}
          hanya terdefinisi benar untuk jumlah satu.

    \vspace{0.22em}
    \item \asm{ROL}/\asm{ROR} memutar bit secara melingkar tanpa
          kehilangan informasi, sedangkan \asm{RCL}/\asm{RCR}
          menyertakan \flag{CF} di ujung lingkaran putar.

    \vspace{0.22em}
    \item Pasangan \asm{SAR} dan \asm{RCR} menggeser bilangan 32 bit
          (DX:AX) sebagai satu kesatuan.

    \vspace{0.22em}
    \item Geser satu bit setara kali atau bagi dua: \asm{SHL} bersama
          \asm{ADD} menggantikan \asm{MUL} yang boros siklus, tetapi
          pembagian bertanda hanya sah dengan \asm{SAR}.
\end{itemize}
\end{frame}

% =========================================================
% BAGIAN 3
% =========================================================
\PemisahBagian{Studi Kasus Manipulasi Bit}

% ---------------------------------------------------------
\begin{frame}{Memeriksa Paritas Sebuah Byte}
\small \term{Paritas} menyatakan apakah cacah bit 1 sebuah byte genap
atau ganjil; pemeriksaan ini mendeteksi satu bit yang berubah saat
transmisi dan dilayani \flag{PF}.

\vspace{0.2em}
\footnotesize
\begin{center}
\begin{tabular}{@{}p{2.0cm}p{3.9cm}p{7.0cm}@{}}
    \toprule
    \textbf{Cara} & \textbf{Instruksi kunci} & \textbf{Catatan} \\
    \midrule
    \eng{Flag} \flag{PF} & \asm{TEST AL, AL} lalu \asm{JP} &
    Paling singkat; hanya untuk byte rendah hasil operasi \\
    Pelipatan \asm{XOR} & \asm{XOR AL, AH} tiga kali &
    Tidak bergantung \flag{PF}; berlaku untuk
    \eng{word}/\eng{dword} \\
    Pergeseran \asm{ADC} & \asm{SHR AL, 1} lalu \asm{ADC BL, 0} &
    Paling gamblang; iterasi sebanyak cacah bit operand \\
    \bottomrule
\end{tabular}
\end{center}

\vspace{0.2em}
{\footnotesize\color{obeGray}Byte uji \bin{1011 0110} memiliki lima bit
1 --- cacah ganjil, sehingga program mencetak \code{GANJIL}.\par}
\Sumber{Bab XI, Memeriksa Paritas Sebuah Byte}
\end{frame}

% ---------------------------------------------------------
\begin{frame}[fragile]{Paritas Tanpa \flag{PF}: Pelipatan \asm{XOR}}
\small Byte dilipat menjadi setengah lebarnya dengan \asm{XOR}, lalu
diulang sampai tersisa satu bit.

\vspace{0.25em}
\begin{kodeasm}[numbers=none]
MOV AL, nilai       ; byte yang diperiksa
MOV AH, AL
SHR AH, 4
XOR AL, AH          ; lipat nibble tinggi ke nibble rendah
MOV AH, AL
SHR AH, 2
XOR AL, AH
MOV AH, AL
SHR AH, 1
XOR AL, AH
AND AL, 1           ; 1 = cacah bit 1 ganjil, 0 = genap
\end{kodeasm}

\vspace{0.25em}
{\footnotesize\color{obeGray}Pelipatan berturut-turut 4, 2, lalu 1 bit
memindahkan seluruh bit ke posisi terendah, sehingga \asm{AND AL, 1}
menyisakan paritasnya.\par}
\Sumber{Bab XI, Memeriksa Paritas Sebuah Byte}
\end{frame}

% ---------------------------------------------------------
\begin{frame}[fragile]{\eng{Packing} dan \eng{Unpacking} Dua Digit BCD}
\small \term{BCD} menyimpan bilangan desimal empat bit per digit: satu
byte menampung puluhan pada nibble tinggi dan satuan pada nibble
rendah.

\vspace{0.15em}
\footnotesize
\begin{center}
\begin{tabular}{@{}p{3.0cm}p{6.0cm}p{3.9cm}@{}}
    \toprule
    \textbf{Operasi} & \textbf{Instruksi} & \textbf{Hasil} \\
    \midrule
    \eng{Packing} & \asm{MOV AL, puluhan} lalu \asm{SHL AL, CL}
    (CL = 4) lalu \asm{OR AL, satuan} & \bin{0111 0101} = \hex{75} \\
    Ambil satuan & \asm{AND AL, 0Fh} & \bin{0000 0101} = \hex{05} \\
    Ambil puluhan & \asm{SHR AL, CL} dengan CL = 4 &
    \bin{0000 0111} = \hex{07} \\
    \bottomrule
\end{tabular}
\end{center}

\vspace{0.15em}
\begin{kodeasm}[numbers=none]
MOV AL, 7           ; digit puluhan
MOV CL, 4
SHL AL, CL          ; AL = 70h (0111 0000b)
OR  AL, 5           ; AL = 75h -> 75 dalam BCD
MOV AH, AL
AND AH, 0Fh         ; AH = 05h, digit satuan
SHR AL, CL          ; AL = 07h, digit puluhan
\end{kodeasm}
\Sumber{Bab XI, Packing dan Unpacking Dua Digit BCD}
\end{frame}

% ---------------------------------------------------------
\begin{frame}[fragile]{\eng{XOR Cipher} Sederhana}
\small Bila $A$ di-XOR dengan kunci $K$, hasilnya $B = A \oplus K$.
Nilai asli dipulihkan dengan meng-XOR-kan $B$ sekali lagi dengan kunci
yang sama:
$B \oplus K = (A \oplus K) \oplus K = A$.

\vspace{0.2em}
\begin{columns}[t]
\begin{column}{0.58\textwidth}
\begin{kodeasm}[numbers=none]
LEA SI, asli        ; SI <- penunjuk sumber
LEA DI, buffer      ; DI <- penunjuk tujuan
ENKRIPSI:
MOV AL, [SI]        ; ambil satu byte
CMP AL, '$'         ; penanda akhir string?
JE  TAMPIL_ENK
XOR AL, kunci       ; XOR dengan kunci 20h
MOV [DI], AL        ; simpan hasil
INC SI
INC DI
JMP ENKRIPSI
\end{kodeasm}
\end{column}
\begin{column}{0.39\textwidth}
{\color{obeTeal}\bfseries Keluaran}\\[0.15em]
\begin{keluaran}
Teks asli   : BELAJARASM
Terenkripsi : belajarasm
Terdekripsi : BELAJARASM
\end{keluaran}
\end{column}
\end{columns}

\vspace{0.15em}
{\footnotesize\color{obeGray}Dekripsi memakai perulangan yang sama,
tetapi menulis kembali di tempat yang sama. Kunci \hex{20} dipilih agar
hasil enkripsi tetap terbaca.\par}
\Sumber{Bab XI, XOR Cipher Sederhana}
\end{frame}

% ---------------------------------------------------------
\begin{frame}{Rangkuman Bagian 3}
\small
\begin{itemize}
    \item Paritas byte ditentukan oleh cacah bit 1: \asm{TEST AL, AL}
          lalu \asm{JP}, pelipatan \asm{XOR} tiga kali, atau
          \asm{SHR} delapan kali dengan \asm{ADC BL, 0}.

    \vspace{0.25em}
    \item Pelipatan \asm{XOR} tidak bergantung pada \flag{PF} dan mudah
          diperluas ke word maupun \eng{dword}.

    \vspace{0.25em}
    \item \eng{Packing} BCD menyusun dua digit dengan \asm{SHL} lalu
          \asm{OR}; \eng{unpacking} memisahkannya dengan \asm{AND 0Fh}
          dan \asm{SHR} empat bit.

    \vspace{0.25em}
    \item \eng{XOR cipher} dapat dibalik karena
          $A \oplus K \oplus K = A$, sehingga kunci yang sama dipakai
          untuk enkripsi dan dekripsi.

    \vspace{0.25em}
    \item Namun cara itu \emph{bukan} pengamanan sungguh-sungguh: bila
          sebagian teks dapat ditebak, kuncinya terhitung dengan satu
          operasi \asm{XOR}.
\end{itemize}
\end{frame}

% =========================================================
% CONTOH PROGRAM
% =========================================================
\begin{frame}[fragile]{Contoh 1: Memeriksa Paritas Byte}
\begin{columns}[t]
\begin{column}{0.57\textwidth}
\begin{kodeasm}[numbers=none]
MOV AL, nilai       ; byte yang diperiksa
MOV BL, 0           ; penghitung cacah bit 1
MOV CX, 8           ; delapan bit diperiksa
ULANG:
SHR AL, 1           ; bit terkanan pindah ke CF
ADC BL, 0           ; BL = BL + CF
LOOP ULANG          ; ulangi sampai CX habis
TEST BL, 1          ; uji bit terendah cacah bit 1
JNZ CETAK_GANJIL    ; ZF = 0 -> cacah ganjil
MOV AH, 09h         ; cacah genap: cetak GENAP
LEA DX, genap
INT 21h
JMP SELESAI
CETAK_GANJIL:       ; cacah ganjil: cetak GANJIL
\end{kodeasm}
\end{column}
\begin{column}{0.41\textwidth}
{\color{obeTeal}\bfseries Keluaran}\\[0.15em]
\begin{keluaran}
Byte 10110110b berparitas: GANJIL
\end{keluaran}

\vspace{0.4em}
{\footnotesize \asm{ADC BL, 0} menjumlahkan penghitung dengan
\flag{CF} --- di sini \eng{carry} dipakai sebagai sumber bit, bukan
sisa penjumlahan.\\[0.25em]
Byte \bin{1011 0111} akan mencetak \code{GENAP}.\\[0.3em]
\color{obeGray}\file{code/asm/paritas.asm}\par}
\end{column}
\end{columns}
\end{frame}

% ---------------------------------------------------------
\begin{frame}[fragile]{Contoh 2: Enkripsi dan Dekripsi \eng{XOR Cipher}}
\begin{columns}[t]
\begin{column}{0.59\textwidth}
\begin{kodeasm}[numbers=none]
LEA SI, asli        ; SI <- penunjuk teks asli
LEA DI, buffer      ; DI <- penunjuk penampung
ENKRIPSI:
MOV AL, [SI]        ; ambil satu byte teks asli
CMP AL, '$'         ; sudah sampai penanda akhir?
JE  TAMPIL_ENK      ; ya -> hentikan perulangan
XOR AL, kunci       ; XOR dengan kunci 20h
MOV [DI], AL        ; simpan byte terenkripsi
INC SI              ; maju ke byte berikutnya
INC DI
JMP ENKRIPSI
; DEKRIPSI memakai perulangan sama, XOR ulang di tempat yang sama
\end{kodeasm}
\end{column}
\begin{column}{0.39\textwidth}
{\color{obeTeal}\bfseries Keluaran}\\[0.15em]
\begin{keluaran}
Teks asli   : BELAJARASM
Terenkripsi : belajarasm
Terdekripsi : BELAJARASM
\end{keluaran}

\vspace{0.35em}
{\footnotesize Perulangan berhenti pada penanda \code{'\$'}, sehingga
panjang string tidak perlu dihitung.\\[0.25em]
\color{obeGray}\file{code/asm/xor-cipher.asm}\par}
\end{column}
\end{columns}
\end{frame}

% ---------------------------------------------------------
\begin{frame}[fragile]{Contoh 3: Memanipulasi Bidang Bit Status}
\small Satu byte menampung beberapa keadaan sekaligus: bit 0 lampu
menyala, bit 3 pintu terbuka, bit 5 alarm aktif.

\vspace{0.2em}
\begin{kodeasm}[numbers=none]
.DATA
    status  DB 00000000b        ; semula seluruh bit padam
.CODE
    OR  status, 00100001b       ; nyalakan lampu (bit 0) + alarm (bit 5)
    AND status, 11011111b       ; padamkan alarm (bit 5)
    XOR status, 00001000b       ; balikkan pintu (bit 3) -> terbuka
    XOR status, 00001000b       ; balikkan lagi -> tertutup
    TEST status, 00001000b      ; uji bit 3, AL tidak berubah
    JZ  PINTU_TERTUTUP          ; ZF = 1 -> pintu tertutup
\end{kodeasm}

\vspace{0.2em}
{\footnotesize\color{obeGray}Jauh lebih hemat daripada satu variabel per
keadaan: satu byte menampung delapan keadaan dan pemeriksaannya cukup
satu \asm{TEST} diikuti lompatan bersyarat.\par}
\Sumber{Bab XI, Contoh 3: Memanipulasi Bidang Bit Status}
\end{frame}

% =========================================================
% PRAKTIKUM
% =========================================================
\begin{frame}{Praktikum: Manipulasi Bit}
\begin{block}{Judul praktikum}
\small Manipulasi bit --- pemeriksaan paritas dan \eng{packing} BCD
\quad\textbullet\quad \textbf{Sub-CPMK:} 3.2 \quad\textbullet\quad
\textbf{Waktu:} 100 menit \quad\textbullet\quad \textbf{Bentuk:}
individu
\end{block}

\vspace{0.1em}
\begin{columns}[t]
\begin{column}{0.48\textwidth}
{\color{obeBlue}\bfseries Alat dan bahan}
\begin{itemize}\small
    \item Komputer bersistem operasi Windows
    \item Emu8086 atau GUI Turbo Assembler + TASM dan TLINK
    \item Berkas \file{paritas.asm} dan \file{xor-cipher.asm}
    \item Lembar pengamatan paritas
\end{itemize}
\end{column}
\begin{column}{0.48\textwidth}
{\color{obeBlue}\bfseries Langkah kerja}
\begin{enumerate}\small
    \item Telusuri \file{paritas.asm} dengan \file{Single Step};
          catat \reg{AL}, \reg{BL}, \flag{CF}.
    \item Uji byte \bin{1011 0111} dan \bin{0101 0101}.
    \item Nyalakan, padamkan, balikkan, lalu uji bit ke-3 \hex{A5}.
    \item Susun BCD $75$, pisahkan kembali, ulangi untuk $93$.
    \item Bandingkan \asm{SHR} dan \asm{SAR} pada \hex{FF96};
          hitung $7 \times 12$ tanpa \asm{MUL}.
\end{enumerate}
\end{column}
\end{columns}

\vspace{0.05em}
{\footnotesize\textbf{Dikumpulkan:} tabel pengamatan paritas, program
manipulasi bit ke-$n$, program \eng{packing} BCD beserta tangkapan
layar, tabel \asm{SHR} vs \asm{SAR}, dan lembar \textbf{Kuis 3}.\par}
\end{frame}

% =========================================================
% LATIHAN
% =========================================================
\begin{frame}{Latihan}
\small
\begin{enumerate}
    \item Jelaskan dengan kalimat Anda sendiri perbedaan cara kerja
          \asm{AND}, \asm{OR}, dan \asm{XOR} pada level bit. Sertakan
          tabel kebenaran singkat untuk masing-masing instruksi.

    \vspace{0.28em}
    \item Mengapa \flag{CF} dan \flag{OF} selalu bernilai 0 setelah
          \asm{AND}, \asm{OR}, \asm{XOR}, atau \asm{TEST}? Kaitkan
          dengan sifat operasi \eng{bitwise} yang tidak menjalarkan
          \eng{carry} antarbit.

    \vspace{0.28em}
    \item Sebutkan tiga perbedaan \asm{TEST} dan \asm{CMP}, lalu
          tentukan yang lebih tepat untuk memeriksa apakah bit ke-4
          \reg{AL} menyala, dan untuk memeriksa apakah \reg{AX} lebih
          besar daripada \reg{BX}.

    \vspace{0.28em}
    \item Diketahui \reg{AL} = \hex{A5}. Tentukan nilai akhir \reg{AL}
          dan nilai \flag{ZF} untuk potongan program berikut:
          \asm{AND AL, 00001111b}; \asm{OR AL, 00110000b};
          \asm{XOR AL, 00000011b}; \asm{TEST AL, AL}.

    \vspace{0.28em}
    \item Diketahui \reg{AL} = \bin{1001 0110} (\hex{96}) yang
          ditafsirkan bertanda. Tentukan hasil dan nilai \flag{CF}
          untuk \asm{SHL AL, 1}, \asm{SHR AL, 1}, dan \asm{SAR AL, 1};
          jelaskan mengapa hanya \asm{SAR} yang memberi hasil
          pembagian bertanda yang benar.
\end{enumerate}
\end{frame}

% =========================================================
% RANGKUMAN
% =========================================================
\begin{frame}{Rangkuman}
\small
\begin{itemize}
    \item \asm{AND}, \asm{OR}, \asm{XOR}, dan \asm{NOT} bekerja pada
          setiap posisi bit secara mandiri tanpa menjalarkan
          \eng{carry}: \asm{AND} untuk \eng{masking}, \asm{OR} untuk
          menyalakan, \asm{XOR} untuk membalik, dan \asm{NOT} membalik
          seluruh bit tanpa menyentuh \eng{flag}.

    \vspace{0.2em}
    \item Setelah \asm{AND}, \asm{OR}, \asm{XOR}, dan \asm{TEST},
          \flag{CF} dan \flag{OF} selalu 0, sedangkan \flag{ZF},
          \flag{SF}, dan \flag{PF} mengikuti hasil.

    \vspace{0.2em}
    \item Manipulasi bit ke-$n$ memakai mask $2^n$: \asm{OR} untuk
          menyalakan, \asm{AND} dengan komplemen mask untuk memadamkan,
          \asm{XOR} untuk membalik, \asm{TEST} untuk menguji.

    \vspace{0.2em}
    \item \asm{TEST} = \asm{AND} tanpa menyimpan hasil (memeriksa bit);
          \asm{CMP} = \asm{SUB} tanpa menyimpan hasil (membandingkan
          nilai utuh).

    \vspace{0.2em}
    \item \asm{SAL} identik dengan \asm{SHL}, sedangkan \asm{SAR}
          menyisipkan salinan bit tanda; \asm{ROL}/\asm{ROR} memutar
          melingkar dan \asm{RCL}/\asm{RCR} menyertakan \flag{CF}.
          Bit yang keluar selalu ditampung \flag{CF}.

    \vspace{0.2em}
    \item Tiga studi kasus: pemeriksaan paritas (\flag{PF}, pelipatan
          \asm{XOR}, atau \asm{SHR} + \asm{ADC}), \eng{packing} BCD,
          dan \eng{XOR cipher} yang dapat dibalik karena
          $A \oplus K \oplus K = A$.
\end{itemize}
\end{frame}

% ---------------------------------------------------------
\begin{frame}{Cek Pemahaman}
\small
\begin{enumerate}
    \item Nilai \flag{CF} dan \flag{OF} setelah instruksi
          \asm{XOR AL, BL} dijalankan adalah \dots
          \begin{enumerate}\footnotesize
              \item keduanya 0
              \item keduanya 1
              \item \flag{CF} = 1 dan \flag{OF} = 0
              \item bergantung pada nilai \reg{AL} dan \reg{BL}
          \end{enumerate}

    \vspace{0.25em}
    \item Instruksi untuk memadamkan bit ke-3 \reg{AL} tanpa mengubah
          bit lain adalah \dots
          \begin{enumerate}\footnotesize
              \item \asm{OR AL, 08h} \quad
              \item \asm{AND AL, 08h} \quad
              \item \asm{AND AL, 0F7h} \quad
              \item \asm{XOR AL, 08h}
          \end{enumerate}

    \vspace{0.25em}
    \item Diketahui \reg{AL} = \bin{1111 1111} yang ditafsirkan
          sebagai $-1$. Instruksi \asm{SHR AL, 1} menghasilkan
          nilai \dots
          \begin{enumerate}\footnotesize
              \item $-1$ karena bit tanda dipertahankan
              \item $+127$ karena bit tanda diisi nol
              \item $0$ karena seluruh bit bergeser keluar
              \item $-64$ karena hasilnya dibulatkan ke bawah
          \end{enumerate}
\end{enumerate}
\end{frame}

% =========================================================
% PENUTUP
% =========================================================
\begin{frame}{Penutup}
\small
\begin{block}{Sebelum pertemuan berikutnya}
\begin{itemize}
    \item Selesaikan praktikum: tabel pengamatan paritas, program
          manipulasi bit ke-$n$, program \eng{packing} BCD, dan tabel
          perbandingan \asm{SHR} dengan \asm{SAR}.
    \item Latih kembali penelusuran bit secara manual pada byte
          \hex{96} dan \hex{A5} sebelum Kuis 3.
    \item Baca \textbf{Bab XII} --- Struktur Kontrol dan Percabangan.
\end{itemize}
\end{block}

\vspace{0.3em}
{\footnotesize\textbf{Bacaan lanjutan:} Bab XI bagian Ringkasan
Kedelapan Instruksi Geser dan Rotasi serta Studi Kasus Manipulasi Bit;
daftar \eng{encoding} dan waktu siklus Intel 8086 (perbandingan
\asm{SHL} dan \asm{MUL}); pembahasan \eng{flag} secara mendalam pada
Bab V.\par}

\vspace{0.35em}
\begin{center}
{\color{obeBlue}\sffamily\large\bfseries Terima kasih\par}
\vspace{0.15em}
{\color{obeGray}\sffamily\small \NamaMK\ (\KodeMK) \textbullet\
\SemesterTeks\par}
\end{center}
\end{frame}

\end{document}
